%% kapa: exam=1 round=23 year=2012 period=1 form=A subject=economics subject_ko=경제원론 seq=3 number=43
%% source: Q-Net 감정평가사 기출문제 (cst00309) artlSeq=5052578 file=raw/1008/5052578/2012년 제23회 감정평가사 제1차시험 기출문제(1교시 A형).pdf
%% conversion: kordoc 4.24.51 -> tools/gen1.py (auto); review=proofread-sonnet; flags=none
\begin{question}{43}
가격차별에 관한 설명으로 옳지 않은 것은?
\begin{choices}
  \item 1급 가격차별을 하면 소비자 잉여는 모두 생산자 잉여가 된다.
  \item 완전경쟁시장과 가격차별은 양립하지 않는다.
  \item 가격차별은 경제적 순손실(deadweight loss)을 항상 증대시킨다.
  \item 가격차별은 독점기업의 이윤극대화 전략 중의 하나이다.
  \item 극장에서의 조조할인 요금제는 가격차별의 일종이다.
\end{choices}
\end{question}
\answer{3}
